\section{M}

% machen (to make)
\begin{verbentry}{
  style = {acc},
  toc = {machen (to make)},
  name = {machen},
  english = {to make},
  type = {regular, + Akkusative},
  auxiliary = {haben}
}
 
    
     \vspace{5pt}

    \hrule
    
    \vspace{5pt}

  \begin{tabular}{rl}
     \multicolumn{2}{l}{\textbf{PRÄSENS} }\\
    ich & \textbf{mache}\\
    du & \textbf{machst}\\
    er/sie/es & \textbf{macht}\\
    wir & \textbf{machen}\\
    ihr & \textbf{macht}\\
    sie/Sie & \textbf{machen}\\
  \end{tabular}
  \hfill
     \begin{tabular}{rl}
    \multicolumn{2}{l}{\textbf{PRÄTERITUM} }\\
    ich & \textbf{machte}\\
    du & \textbf{machtest}\\
    er/sie/es & \textbf{machte}\\
    wir & \textbf{machten}\\
    ihr & \textbf{machtet}\\
    sie/Sie & \textbf{machten}\\
  \end{tabular}
  \hfill
  \begin{tabular}{rl}
     \multicolumn{2}{l}{\textbf{IMPERATIV} }\\
   & \textbf{mach(e) (du)!}\\
   & \textbf{machen wir!}\\
   & \textbf{macht (ihr)!}\\
   & \textbf{machen Sie!}\\
   & \vspace{10pt}
  \end{tabular}

    \vspace{10pt}

 \begin{tabular}{rl}
     \multicolumn{2}{l}{\textbf{PERFEKT}}\\
   & haben + \textbf{gemacht}\\
  \end{tabular}
\hfill
   \begin{tabular}{rl}
   
    \multicolumn{2}{l}{\textbf{FUTURE I} }\\
   & werden + \textbf{machen}\\
 
  \end{tabular}
\hfill
   \begin{tabular}{rl}
      \multicolumn{2}{l}{\textbf{PLUSQUAMPERFEKT}}\\
   & hatten + \textbf{gemacht}\\
  \end{tabular}
  

  \vspace{10pt}

   \begin{tabular}{lll}
      \multicolumn{3}{l}{\textbf{MIT PRÄPOSITIONEN}}\\
& \textbf{machen mit} + Dat & (to participate in)\\
& \textbf{machen für} + Akk & (to do something for someone / purpose)\\
&\textbf{machen aus} + Dat & (to turn something into something)\\
&\textbf{machen über} + Akk & (to complain / joke about)\\
&\textbf{machen von} + Dat & (to make of / think of)\\
&\textbf{machen bei} + Dat & (to take part in)\\
\vspace{40pt}
  \end{tabular}
  \hfill
   \begin{tabular}{lll}
      \multicolumn{3}{l}{\textbf{TRENNBARE VERBEN}}\\
 & \textbf{an$|$machen} & (to turn on)\\
 & \textbf{aus$|$machen} & (to turn off)\\
 & \textbf{auf$|$machen} & (to open)\\
  & \textbf{zu$|$machen} & (to close)\\
  & \textbf{mit$|$machen} & (to participate)\\
  & \textbf{nach$|$machen} & (to imitate)\\
  & \textbf{weiter$|$machen} & (to continue)\\
     & \textbf{fertig$|$machen} & (to finish /exhaust)\\
     & \textbf{kaputt$|$machen} & (to break)\\
     & \textbf{los$|$machen} & (to loosen / detach)\\
  \end{tabular}


  \vspace{-5pt}
\end{verbentry}

% mauern (to build walls / brick)
\begin{verbentry}{
  style = {default},
  toc = {mauern (to build walls / brick)},
  name = {mauern},
  english = {to build walls / brick},
  type = {regular},
  auxiliary = {haben}
}
 
    
     \vspace{5pt}

    \hrule
    
    \vspace{5pt}

  \begin{tabular}{rl}
     \multicolumn{2}{l}{\textbf{PRÄSENS} }\\
    ich & \textbf{mauere}\\
    du & \textbf{mauerst}\\
    er/sie/es & \textbf{mauert}\\
    wir & \textbf{mauern}\\
    ihr & \textbf{mauert}\\
    sie/Sie & \textbf{mauern}\\
  \end{tabular}
  \hfill
  \begin{tabular}{rl}
    \multicolumn{2}{l}{\textbf{PRÄTERITUM} }\\
    ich & \textbf{mauerte}\\
    du & \textbf{mauertest}\\
    er/sie/es & \textbf{mauerte}\\
    wir & \textbf{mauerten}\\
    ihr & \textbf{mauertet}\\
    sie/Sie & \textbf{mauerten}\\
  \end{tabular}
  \hfill
  \begin{tabular}{rl}
     \multicolumn{2}{l}{\textbf{IMPERATIV} }\\
   & \textbf{mauere (du)!}\\
   & \textbf{mauern wir!}\\
   & \textbf{mauert (ihr)!}\\
   & \textbf{mauern Sie!}\\
   & \vspace{10pt}
  \end{tabular}

 \vspace{10pt}

 \begin{tabular}{rl}
     \multicolumn{2}{l}{\textbf{PERFEKT}}\\
   & haben + \textbf{gemauert}\\
  \end{tabular}
\hfill
\begin{tabular}{rl}
    \multicolumn{2}{l}{\textbf{FUTURE I}}\\
   & werden + \textbf{mauern}\\
  \end{tabular}
\hfill
\begin{tabular}{rl}
    \multicolumn{2}{l}{\textbf{PLUSQUAMPERFEKT}}\\
   & hatten + \textbf{gemauert}\\
  \end{tabular}

 \vspace{10pt}

\begin{tabular}{lll}
      \multicolumn{3}{l}{\textbf{MIT PRÄPOSITIONEN}}\\
& \textbf{mauern mit} + Dat & (to build with / using)\\
& \textbf{mauern an} + Dat & (to build onto / on)\\
\end{tabular}
  \hfill
\begin{tabular}{lll}
      \multicolumn{3}{l}{\textbf{TRENNBARE VERBEN}}\\
 & \textbf{ein$|$mauern} & (to wall in / enclose)\\
\vspace{10pt}
  \end{tabular}

  \vspace{-5pt}
\end{verbentry}

% meinen (to mean something)
\begin{verbentry}{
  style = {acc},
  toc = {meinen (to mean something)},
  name = {meinen},
  english = {to mean something},
  type = {regular, + Akkusativ},
  auxiliary = {haben}
}
 
    
     \vspace{5pt}

    \hrule
    
    \vspace{5pt}

  \begin{tabular}{rl}
     \multicolumn{2}{l}{\textbf{PRÄSENS} }\\
    ich & \textbf{meine}\\
    du & \textbf{meinst}\\
    er/sie/es & \textbf{meint}\\
    wir & \textbf{meinen}\\
    ihr & \textbf{meint}\\
    sie/Sie & \textbf{meinen}\\
  \end{tabular}
  \hfill
     \begin{tabular}{rl}
    \multicolumn{2}{l}{\textbf{PRÄTERITUM} }\\
    ich & \textbf{meinte}\\
    du & \textbf{meintest}\\
    er/sie/es & \textbf{meinte}\\
    wir & \textbf{meinten}\\
    ihr & \textbf{meintet}\\
    sie/Sie & \textbf{meinten}\\
  \end{tabular}
  \hfill
  \begin{tabular}{rl}
     \multicolumn{2}{l}{\textbf{IMPERATIV} }\\
   & \textbf{mein(e) (du)!}\\
   & \textbf{meinen wir!}\\
   & \textbf{meint (ihr)!}\\
   & \textbf{meinen Sie!}\\
   & \vspace{10pt}
  \end{tabular}

    \vspace{10pt}

 \begin{tabular}{rl}
     \multicolumn{2}{l}{\textbf{PERFEKT}}\\
  &  haben + \textbf{gemeint}\\
  \end{tabular}
\hfill
   \begin{tabular}{rl}
   
    \multicolumn{2}{l}{\textbf{FUTURE I} }\\
   & werden + \textbf{meinen}\\
 
  \end{tabular}
\hfill
   \begin{tabular}{rl}
      \multicolumn{2}{l}{\textbf{PLUSQUAMPERFEKT}}\\
   & hatten + \textbf{gemeint}\\
  \end{tabular}
  
  
\vspace{10pt}

\begin{tabular}{lll}
\multicolumn{3}{l}{\textbf{MIT PRÄPOSITIONEN}}\\
& \textbf{---} & \\
\end{tabular}
\hfill
\begin{tabular}{lll}
\multicolumn{3}{l}{\textbf{TRENNBARE VERBEN}}\\
& \textbf{---} & \\
\end{tabular}
\vspace{-5pt}
\end{verbentry}

% melken (to milk)
\begin{verbentry}{
  style = {default},
  toc = {melken (to milk)},
  name = {melken},
  english = {to milk},
  type = {irregular},
  auxiliary = {haben}
}
 
    
     \vspace{5pt}

    \hrule
    
    \vspace{5pt}

  \begin{tabular}{rl}
     \multicolumn{2}{l}{\textbf{PRÄSENS} }\\
    ich & \textbf{melke}\\
    du & \textbf{melkst}\\
    er/sie/es & \textbf{melkt}\\
    wir & \textbf{melken}\\
    ihr & \textbf{melkt}\\
    sie/Sie & \textbf{melken}\\
  \end{tabular}
  \hfill
     \begin{tabular}{rl}
    \multicolumn{2}{l}{\textbf{PRÄTERITUM} }\\
    ich & \textbf{molk}\\
    du & \textbf{molkst}\\
    er/sie/es & \textbf{molk}\\
    wir & \textbf{molken}\\
    ihr & \textbf{molkt}\\
    sie/Sie & \textbf{molken}\\
  \end{tabular}
  \hfill
  \begin{tabular}{rl}
     \multicolumn{2}{l}{\textbf{IMPERATIV} }\\
   & \textbf{melk(e) (du)!}\\
   & \textbf{melken wir!}\\
   & \textbf{melkt (ihr)!}\\
   & \textbf{melken Sie!}\\
   & \vspace{10pt}
  \end{tabular}

    \vspace{10pt}

 \begin{tabular}{rl}
     \multicolumn{2}{l}{\textbf{PERFEKT}}\\
   & haben + \textbf{gemolken}\\
  \end{tabular}
\hfill
   \begin{tabular}{rl}
   
    \multicolumn{2}{l}{\textbf{FUTURE I} }\\
   & werden + \textbf{melken}\\
 
  \end{tabular}
\hfill
   \begin{tabular}{rl}
      \multicolumn{2}{l}{\textbf{PLUSQUAMPERFEKT}}\\
   & hatten + \textbf{gemolken}\\
  \end{tabular}

   \vspace{10pt}

   \begin{tabular}{lll}
      \multicolumn{3}{l}{\textbf{MIT PRÄPOSITIONEN}}\\
& \textbf{melken von} + Dat & (to extract from)\\
& \textbf{melken aus} + Dat & (to extract out of)\\
   \end{tabular}
  \hfill
   \begin{tabular}{lll}
      \multicolumn{3}{l}{\textbf{TRENNBARE VERBEN}}\\
 & \textbf{ab$|$melken} & (to siphon off / exploit)\\
 & \textbf{aus$|$melken} & (to milk out / exploit fully)\\
  \end{tabular}

\vspace{-5pt}
\end{verbentry}

% merken (to notice / remember)
\begin{verbentry}{
  style = {default},
  toc = {merken (to notice / remember)},
  name = {merken},
  english = {to notice / remember},
  type = {regular},
  auxiliary = {haben}
}


\vspace{5pt}
\hrule
\vspace{5pt}

\begin{tabular}{rl}
\multicolumn{2}{l}{\textbf{PRÄSENS}}\\
ich & \textbf{merke}\\
du & \textbf{merkst}\\
er/sie/es & \textbf{merkt}\\
wir & \textbf{merken}\\
ihr & \textbf{merkt}\\
sie/Sie & \textbf{merken}\\
\end{tabular}
\hfill
\begin{tabular}{rl}
\multicolumn{2}{l}{\textbf{PRÄTERITUM}}\\
ich & \textbf{merkte}\\
du & \textbf{merktest}\\
er/sie/es & \textbf{merkte}\\
wir & \textbf{merkten}\\
ihr & \textbf{merktet}\\
sie/Sie & \textbf{merkten}\\
\end{tabular}
\hfill
\begin{tabular}{rl}
\multicolumn{2}{l}{\textbf{IMPERATIV}}\\
& \textbf{merk(e) (du)!}\\
& \textbf{merken wir!}\\
& \textbf{merkt (ihr)!}\\
& \textbf{merken Sie!}\\
\end{tabular}

\vspace{10pt}

\begin{tabular}{rl}
\multicolumn{2}{l}{\textbf{PERFEKT}}\\
& haben + \textbf{gemerkt}\\
\end{tabular}
\hfill
\begin{tabular}{rl}
\multicolumn{2}{l}{\textbf{FUTURE I}}\\
& werden + \textbf{merken}\\
\end{tabular}
\hfill
\begin{tabular}{rl}
\multicolumn{2}{l}{\textbf{PLUSQUAMPERFEKT}}\\
& hatten + \textbf{gemerkt}\\
\end{tabular}

\vspace{10pt}

\begin{tabular}{lll}
\multicolumn{3}{l}{\textbf{MIT PRÄPOSITIONEN}}\\
& \textbf{merken an} + Dat & (to notice in)\\
& \textbf{merken sich} + Akk & (to remember)\\
\end{tabular}
\hfill
\begin{tabular}{lll}
\multicolumn{3}{l}{\textbf{TRENNBARE VERBEN}}\\
& \textbf{fest$|$merken} & (to take note of)\\
& \textbf{zurück$|$merken} & (to remember for later)\\
\end{tabular}
\vspace{-5pt}
\end{verbentry}

% missverstehen (to misunderstand)
\begin{verbentry}{
  style = {default},
  toc = {missverstehen (to misunderstand)},
  name = {missverstehen},
  english = {to misunderstand},
  type = {irregular},
  auxiliary = {haben}
}
 
    
     \vspace{5pt}

    \hrule
    
    \vspace{5pt}

  \begin{tabular}{rl}
     \multicolumn{2}{l}{\textbf{PRÄSENS} }\\
    ich & \textbf{missverstehe}\\
    du & \textbf{missverstehst}\\
    er/sie/es & \textbf{missversteht}\\
    wir & \textbf{missverstehen}\\
    ihr & \textbf{missversteht}\\
    sie/Sie & \textbf{missverstehen}\\
  \end{tabular}
  \hfill
     \begin{tabular}{rl}
    \multicolumn{2}{l}{\textbf{PRÄTERITUM} }\\
    ich & \textbf{missverstand}\\
    du & \textbf{missverstandest}\\
    er/sie/es & \textbf{missverstand}\\
    wir & \textbf{missverstanden}\\
    ihr & \textbf{missverstandet}\\
    sie/Sie & \textbf{missverstanden}\\
  \end{tabular}
  \hfill
  \begin{tabular}{rl}
     \multicolumn{2}{l}{\textbf{IMPERATIV} }\\
   & \textbf{missversteh(e) (du)!}\\
   & \textbf{missverstehen wir!}\\
   & \textbf{missversteht (ihr)!}\\
   & \textbf{missverstehen Sie!}\\
   & \vspace{10pt}
  \end{tabular}

    \vspace{10pt}

 \begin{tabular}{rl}
     \multicolumn{2}{l}{\textbf{PERFEKT}}\\
   & haben + \textbf{missverstanden}\\
  \end{tabular}
\hfill
   \begin{tabular}{rl}
    \multicolumn{2}{l}{\textbf{FUTURE I} }\\
   & werde + \textbf{missverstehen}\\
  \end{tabular}
\hfill
   \begin{tabular}{rl}
      \multicolumn{2}{l}{\textbf{PLUSQUAMPERFEKT}}\\
   & hatten + \textbf{missverstanden}\\
  \end{tabular}

   \vspace{10pt}

   \begin{tabular}{lll}
      \multicolumn{3}{l}{\textbf{MIT PRÄPOSITIONEN}}\\
& \textbf{missverstehen als} + Akk & (to misunderstand as)\\
& \textbf{missverstehen wegen} + Gen & (to misunderstand because of)\\
\end{tabular}
  \hfill
   \begin{tabular}{lll}
\multicolumn{3}{l}{\textbf{TRENNBARE VERBEN}}\\
& \textbf{---} & \\
\end{tabular}

  \vspace{-5pt}
\end{verbentry}
